\subsection{Quadratic overlap estimates for ordered blocks}

The following estimates extend the second-order comparison to blocks of
unequal lengths. They refine the first-order error in
\cite[Lemma~1$'$(i), eq.~(S11)]{Malz2024Preparation}; the exponent obtained
here is a project result.

\begin{theorem}[Distance between unit vectors]
    \label{thm:ldp_unit_vector_deficit}
    \lean{one_sub_re_inner_eq_norm_sub_sq_div_two}
    \leanok
    Let $x,y$ be unit vectors in a real or complex inner-product space. Then
    $1-\operatorname{Re}\langle x,y\rangle=\lVert x-y\rVert^2/2$.
\end{theorem}
\begin{proof}\leanok
    Expand $\lVert x-y\rVert^2=\lVert x\rVert^2+\lVert y\rVert^2
        -2\operatorname{Re}\langle x,y\rangle$ and use the unit norms.
\end{proof}

\begin{theorem}[Nonnegative polar compression coefficient]
    \label{thm:ldp_polar_compression_nonnegative}
    \lean{MPSTensor.mixedMapLM_polarPosTensor_fixedPointTensor_trace_nonneg}
    \leanok
    Let $B$ be a tensor, $P$ its positive polar tensor, and $\sigma\ge0$ a
    matrix on its virtual space. Let $P_\infty$ be the fixed-point tensor
    associated with $\sigma$, and put
    $\alpha=\Tr\sum_iP^i\sigma(P_\infty^i)^\dagger$.
    Then $\alpha$ is real and nonnegative.
\end{theorem}
\begin{proof}\leanok
    \uses{lem:ldp_overlapping_overlaps_second_order}
    Write $G$ for the positive polar matrix of the physical matrix of $B$.
    The trace identity in Lemma~\ref{lem:ldp_overlapping_overlaps_second_order}
    gives $\alpha=\Tr(G((\sqrt\sigma)^{\mathsf T}\otimes\sigma))$.
    Both factors are positive semidefinite, so their trace product is
    real and nonnegative.
\end{proof}

\begin{theorem}[Quadratic polar compression deficit]
    \label{thm:ldp_polar_compression_deficit}
    \lean{MPSTensor.mixedMapLM_polarPosTensor_fixedPointTensor_trace_geometry}
    \leanok
    In the setting of Theorem~\ref{thm:ldp_polar_compression_nonnegative},
    suppose $\Tr\sigma=\Tr\E_B(\sigma)=1$. Put
    $\iota(X)=(X^i\sqrt\sigma)_i$, with the norm induced by the coordinate inner product.
    Then $\operatorname{Re}\alpha\le1$ and
    $1-\operatorname{Re}\alpha
        =\lVert\iota(P)-\iota(P_\infty)\rVert^2/2$.
\end{theorem}
\begin{proof}\leanok
    \uses{thm:ldp_unit_vector_deficit,
        thm:ldp_positive_part_overlap_second_order,
        lem:ldp_overlapping_overlaps_second_order,
        thm:ldpolar_blocked_transfer}
    The positive polar tensor has the same transfer map as $B$. Thus
    $\lVert\iota(P)\rVert=\lVert\iota(P_\infty)\rVert=1$ and
    $\langle\iota(P_\infty),\iota(P)\rangle=\alpha$.
    Apply Theorem~\ref{thm:ldp_unit_vector_deficit}; the nonnegativity of the
    squared distance also gives $\operatorname{Re}\alpha\le1$.
\end{proof}

\begin{theorem}[Ordered perturbations with vanishing compression]
    \label{thm:ldp_ordered_idempotent_quadratic}
    \lean{IsIdempotentElem.norm_trace_prod_add_sub_le_of_le}
    \leanok
    Let $e$ be an idempotent in a normed ring and put $f=1-e$. Let $\tau$
    be an additive complex-valued map satisfying $\tau(ab)=\tau(ba)$ and
    $|\tau(a)|\le K\lVert a\rVert$, where $K\ge0$. Suppose that a sequence
    $Z_j$ satisfies $eZ_je=0$ and
    $\lVert fZ_je\rVert,\lVert eZ_jf\rVert,\lVert fZ_jf\rVert\le z\le1/2$.
    For every $M\ge2$,
    \begin{align}
        \left|\tau\left(\prod_{j=0}^{M-1}(e+Z_j)\right)-\tau(e)\right|
         & \le K(\lVert e\rVert+3\lVert f\rVert)
        (2Mz^2)e^{2Mz^2}.
        \label{eq:ldp_ordered_quadratic_trace}
    \end{align}
    All products are ordered with increasing index from left to right.
\end{theorem}
\begin{proof}\leanok
    Put $P_m=\prod_{j=0}^{m-1}(e+Z_j)$ and $W_j=fZ_jf$. For $a=e,f$,
    the blocks $x_m=aP_me$ and $y_m=aP_mf$ satisfy
    $x_{m+1}=x_m+y_mfZ_me$ and
    $y_{m+1}=x_meZ_mf+y_mW_m$.
    With $\rho=1+z^2/(1-z)$, induction gives
    $\lVert eP_me-e\rVert\le\lVert e\rVert(\rho^m-1)$ and
    $\lVert eP_mf\rVert\le\lVert e\rVert z\rho^m/(1-z)$.
    The analogous recurrences starting from $f$ give
    \begin{align}
        \lVert fP_me\rVert
         & \le\lVert f\rVert\frac{z}{1-z}(\rho^m-z^m), \notag \\
        \left\lVert fP_mf-f\prod_{j=0}^{m-1}W_j\right\rVert
         & \le\lVert f\rVert\frac{z^2}{(1-z)^2}(\rho^m-z^m).
        \notag
    \end{align}
    Cyclicity gives $\tau(P_M)=\tau(eP_Me)+\tau(fP_Mf)$.
    Since $M\ge2$, the complementary product has norm at most
    $\lVert f\rVert z^M\le\lVert f\rVert z^2$. Now use
    $\rho\le1+2z^2$ and $(1+2z^2)^M-1\le2Mz^2e^{2Mz^2}$.
\end{proof}

The restriction $M\ge2$ concerns the absolute trace. Indeed, for
$e=\operatorname{diag}(1,0)$ and $Z=\operatorname{diag}(0,t)$,
one has $eZe=0$ but $\Tr(e+Z)-\Tr e=t$. The normalized one-block
preparation error is estimated separately by a squared-distance inequality.

\begin{theorem}[Ordered factors with scalar compression]
    \label{thm:ldp_ordered_compression_quadratic}
    \lean{IsIdempotentElem.exists_norm_trace_prod_sub_one_le_sq}
    \leanok
    Let $\mathcal A$ be a normed complex algebra, let $e^2=e$, and let
    $\tau:\mathcal A\to\mathbb C$ be linear with $\tau(e)=1$,
    $\tau(ab)=\tau(ba)$ and $|\tau(a)|\le K\lVert a\rVert$ for $K\ge0$.
    There is $C>0$, depending only on these data, with the following
    property. For any sequences $T_j\in\mathcal A$, $\alpha_j\in\mathbb C$,
    real $\delta$, and integer $M\ge2$, suppose
    \begin{align}
        eT_je      & =\alpha_je, & \lVert T_j-e\rVert & \le\delta, \notag \\
        |\alpha_j| & \le1,       & |1-\alpha_j|       & \le\delta^2.
        \label{eq:ldp_ordered_scalar_compression}
    \end{align}
    If $CM\delta^2<1$, then
    $\left|\tau(\prod_{j=0}^{M-1}T_j)-1\right|
        \le CM\delta^2e^{CM\delta^2}$.
\end{theorem}
\begin{proof}\leanok
    \uses{thm:ldp_ordered_idempotent_quadratic}
    Choose $C\ge4$ large enough in terms of $K$, $\lVert e\rVert$ and
    $\lVert1-e\rVert$. The smallness condition gives
    $\delta\le1$ and $|\alpha_j|\ge3/4$. Set
    $Z_j=\alpha_j^{-1}T_j-e$. Then $eZ_je=0$ and
    $\lVert Z_j\rVert\le2(1+\lVert e\rVert)\delta$.
    Its three complementary blocks are bounded by a fixed multiple of
    $\delta$, so Theorem~\ref{thm:ldp_ordered_idempotent_quadratic} applies.
    Put $a=\prod_{j=0}^{M-1}\alpha_j$ and
    $P=\prod_{j=0}^{M-1}(e+Z_j)$. Then $|a|\le1$,
    $|a-1|\le M\delta^2e^{M\delta^2}$, and
    $\tau(\prod_{j=0}^{M-1}T_j)-1=a(\tau(P)-1)+(a-1)$.
    Enlarging $C$ gives the assertion.
\end{proof}

\begin{theorem}[Telescoping an ordered product]
    \label{thm:ldp_ordered_telescoping_identity}
    \lean{prod_range_sub_pow_eq_sum_mul_sub_mul}
    \leanok
    In any ring, for a sequence $X_j$, an element $e$, and $M\ge0$,
    \begin{align}
        \prod_{j=0}^{M-1}X_j-e^M
         & =\sum_{j=0}^{M-1}
        \left(\prod_{i=0}^{j-1}X_i\right)(X_j-e)e^{M-1-j}.
        \label{eq:ldp_ordered_telescoping_identity}
    \end{align}
\end{theorem}
\begin{proof}\leanok
    Induct on $M$, separating the last factor in the ordered product and
    the last summand.
\end{proof}

\begin{theorem}[First-order bound near an idempotent]
    \label{thm:ldp_ordered_telescoping_bound}
    \lean{norm_prod_range_sub_pow_le_of_isIdempotentElem}
    \leanok
    In a normed ring, let $e^2=e$, $\lVert e\rVert,\lVert1\rVert\le c$, and
    $\lVert X_j-e\rVert\le\delta$ for every $j$. Then for every $M\ge0$,
    $\lVert\prod_{j=0}^{M-1}X_j-e^M\rVert
        \le c((1+c\delta)^M-1)$.
\end{theorem}
\begin{proof}\leanok
    \uses{thm:ldp_ordered_telescoping_identity}
    Every power of $e$ has norm at most $c$. Apply
    (\ref{eq:ldp_ordered_telescoping_identity}) and induct on $M$, using
    $\lVert\prod_{i=0}^{j-1}X_i\rVert
        \le c+c((1+c\delta)^j-1)$ in each summand.
\end{proof}

\begin{theorem}[Normalization comparison]
    \label{thm:ldp_normalization_triangle}
    \lean{MPSTensor.one_sub_norm_inner_smul_inv_norm_le}
    \leanok
    Let $\psi,v$ lie in a complex inner-product space with
    $\lVert\psi\rVert\le1$. Suppose
    $|1-|\langle\psi,v\rangle||\le a$ and
    $|\lVert v\rVert^2-1|\le b$.
    Define $\widehat v=v/\lVert v\rVert$ when $v\ne0$ and
    $\widehat v=0$ when $v=0$. Then
    $1-|\langle\psi,\widehat v\rangle|\le a+b$.
\end{theorem}
\begin{proof}\leanok
    For $v=0$, the second assumption gives $b\ge1$.
    Otherwise put $c=\lVert v\rVert>0$ and
    $w=|\langle\psi,\widehat v\rangle|\le1$.
    Since $cw=|\langle\psi,v\rangle|$,
    $1-w=1-|\langle\psi,v\rangle|+(c-1)w$.
    Use $(c-1)w\le|c-1|\le|c^2-1|$.
\end{proof}

\begin{theorem}[First-order absolute trace for unequal blocks]
    \label{thm:ldp_unequal_absolute_trace_first_order}
    \lean{MPSTensor.exists_norm_trace_prod_transferMatrix_sub_one_le}
    \leanok
    Let $A$ be normal in the gauge~(\ref{eq:ldp_normal_gauge}) with
    correlation length $\xi$, and let $0<\gamma<1$.
    Write $\tau_\ell$ for the mixed transfer matrix of the positive polar
    tensor of the $\ell$-site block against the fixed-point tensor.
    There is $C>0$, depending only on the tensor, its normalized fixed
    state and the chosen spectral bound and rate, such that for every
    $M\ge1$ and every family $\ell_j\ge q$,
    $|\Tr(\tau_{\ell_0}\cdots\tau_{\ell_{M-1}})-1|
        \le Cy e^{Cy}$, where $y=Me^{-\gamma q/\xi}$.
\end{theorem}
\begin{proof}\leanok
    \uses{lem:ldp_positive_part_overlap, thm:ldp_ordered_telescoping_bound}
    Each mixed transfer matrix is within $Ke^{-\gamma\ell/\xi}$ of the
    fixed-point idempotent $e$, whose positive powers have trace one.
    Since $\ell_j\ge q$, apply
    Theorem~\ref{thm:ldp_ordered_telescoping_bound} with
    $\delta=Ke^{-\gamma q/\xi}$ and then the boundedness of the trace.
\end{proof}

\begin{theorem}[Quadratic absolute trace for unequal blocks]
    \label{thm:ldp_unequal_absolute_trace_second_order}
    \lean{MPSTensor.exists_norm_trace_prod_transferMatrix_sub_one_le_sq}
    \leanok
    In the same gauge, let $0<\gamma<1$. There is $C>0$, depending only
    on the tensor, its normalized fixed state and the chosen spectral
    bound and rate, such that for every $M\ge2$ and every family
    $\ell_j\ge q$, with $u=Me^{-2\gamma q/\xi}$ and $Cu<1$,
    $|\Tr(\tau_{\ell_0}\cdots\tau_{\ell_{M-1}})-1|
        \le Cu e^{Cu}$.
\end{theorem}
\begin{proof}\leanok
    \uses{thm:ldp_ordered_compression_quadratic,
        thm:ldp_polar_compression_nonnegative,
        thm:ldp_polar_compression_deficit,
        thm:ldp_positive_part_overlap_second_order}
    The fixed-point transfer matrix $e$ satisfies $e^2=e$ and
    $\Tr e=1$. Each factor has scalar compression
    $e\tau_\ell e=\alpha_\ell e$.
    Theorems~\ref{thm:ldp_polar_compression_nonnegative} and
    \ref{thm:ldp_polar_compression_deficit} give
    $0\le\alpha_\ell\le1$ and
    $1-\alpha_\ell=\lVert\iota(P_\ell)-\iota(P_\infty)\rVert^2/2$.
    Both the mixed-transfer difference and the weighted tensor
    difference are bounded by $Ke^{-\gamma q/\xi}$, uniformly in
    $\ell\ge q$. Apply
    Theorem~\ref{thm:ldp_ordered_compression_quadratic} and absorb
    $K^2$ into $C$.
\end{proof}
